

\section{三篇文献的 fence 目标对比}

\begin{table}[htbp]
\centering
\caption{三篇文献的 fence 目标对比}
\begin{tabularx}{\textwidth}{|l|X|X|X|}
\hline
 & [1] Hu et al. 2023 & [2] Pan \& Chen 2024 & [3] Wen et al. 2026 \\
\hline
\textbf{控制目标} & label-free 刚性编队合围：目标被围入凸包，且智能体间维持刚性编队（稳态相对位置固定） & 连通性保持下的合围：目标被围入凸包，同时保证通信网络连通性不丧失 & 无预设编队合围：目标始终在凸包内，不要求固定编队构型 \\
\hline
\textbf{目标运动} & 时变速度（建模为外系统，$s_1 \leq 0$；$s_1=0$ 时退化为匀速） & 线性动态 $\ddot{q}_0 = S_{01}q_0 + S_{02}\dot{q}_0$（可生成多项式/正弦轨迹，含匀速、时变速度为特例） & 时变速度（外系统矩阵 $S$ 的特征值不全为纯虚数即可，条件更宽松） \\
\hline
\textbf{目标外系统矩阵} & $S=\begin{bmatrix} 0 & 1 \\ s_1 & 0 \end{bmatrix}$（$s_1 \leq 0$） & $S_0=\begin{bmatrix} 0 & I \\ S_{01} & S_{02} \end{bmatrix}$（$S_{01},S_{02}$ 为已知常数矩阵） & $S=\begin{bmatrix} 0 & 1 \\ -s_1 & -s_2 \end{bmatrix}$（含阻尼项 $s_2$，特征值不要求纯虚数） \\
\hline
\textbf{系统维度} & 二维 & $n$ 维（任意维度） & 二维 \\
\hline
\textbf{额外约束} & 碰撞避免 + label-free（无预设编号） & 碰撞避免 + 连通性保持 + 外部扰动抑制 & 碰撞避免 + 无预设编队 \\
\hline
\textbf{测量信息} & 自身位置 $x_i$、自身速度 $v_i$、与目标的相对位置 $x_i - x_d$、与邻居的相对位置 $x_{ik}$ & 自身位置 $q_i$、与邻居的相对位置 $q_i - q_j$，至少一个智能体可获取目标位置和速度 $x_0=\mathrm{col}(q_0,p_0)$（目标状态通过分布式观测器估计） & 与目标的相对位置 $e_i = x_i - x_0$、与邻居的相对位置 $x_{ij}$ \\
\hline
\end{tabularx}
\end{table}

\textbf{关键差异总结}：
\begin{itemize}
  \item [1] 使用智能体自身的全状态（位置和速度）以及目标位置，通过内模原理估计目标速度，是最“信息密集”的设计。
  \item [2] 仅需位置反馈，通过三个分布式观测器分别估计车辆状态、扰动和目标状态；同时设计势函数保证连通性保持和碰撞避免，对通信拓扑有连通性保持的额外要求。
  \item [3] 仅需相对位置信息（智能体间及智能体-目标间），通过输出反馈观测器估计未测量的状态，信息需求最低。
\end{itemize}